\documentclass[11pt,a4paper]{article}
\usepackage[margin=22mm]{geometry}
\usepackage{amsmath,amssymb,bm,booktabs,array,longtable,xcolor,hyperref,microtype}
\usepackage[T1]{fontenc}
\hypersetup{colorlinks=true,linkcolor=blue!55!black,urlcolor=blue!55!black}
\definecolor{note}{RGB}{238,246,252}
\newcommand{\code}[1]{\texttt{#1}}
\newcommand{\sgnq}{s_{ij}}
\title{Yukawa Potential, Reduced Units, and System Setup\\
\large Paper definition $\rightarrow$ Fortran implementation $\rightarrow$ $N=1000$ inputs}
\author{Working note for the MC--NPT reconstruction}
\date{8 July 2026}

\begin{document}
\maketitle
\begin{abstract}
This note derives the hard-core signed Yukawa model used by Sanz et al.\ (PRL
99, 055501, 2007), maps it to the present Fortran code, and gives the complete
conversion between packing fraction, number density, volume, and box length.
It also records which input values matter for a fresh configuration and which
are overridden by a restart.  The aim is to make every dimensionless number in
the simulation auditable.
\end{abstract}
\tableofcontents

\section{Physical model in the paper}
The paper studies an equimolar binary mixture of monodisperse spheres.  Every
particle has diameter $\sigma$ and charge magnitude $Ze$; the two species have
opposite charge signs.  Define
\begin{equation}
  \sgnq = \begin{cases}
  +1,&\text{equal charges (types 1--1 or 2--2)},\\
  -1,&\text{opposite charges (type 1--2)}.
  \end{cases}
\end{equation}
The screened Coulomb pair potential is
\begin{equation}
\frac{u_{ij}(r)}{k_{\mathrm B}T_{\rm phys}}=
\begin{cases}
+\infty, & r<\sigma,\\[2mm]
\displaystyle \sgnq\,
\frac{Z^2\lambda_{\mathrm B}}
{\sigma(1+\kappa\sigma/2)^2}
\frac{\exp[-\kappa(r-\sigma)]}{r/\sigma},
& \sigma\le r<r_c,\\[3mm]
0,&r\ge r_c.
\end{cases}
\label{eq:paper-potential}
\end{equation}
Here $\lambda_{\mathrm B}=e^2/(\epsilon_s k_{\mathrm B}T_{\rm phys})$ is
the Bjerrum length and $\kappa$ is the inverse Debye screening length.  The
paper writes $\kappa=\sqrt{8\pi\lambda_{\mathrm B}\rho_{\rm salt}}$ for its
electrolyte convention.  The simulation parameters are
\begin{equation}
 \kappa\sigma=6, \qquad r_c=3.5\sigma.
\end{equation}
The exponential is referenced to contact, $r=\sigma$.  Consequently the
magnitude of the finite interaction at contact is especially simple.

\section{Energy scale and reduced variables}
Define the Yukawa contact-energy scale
\begin{equation}
 \epsilon_Y \equiv k_{\mathrm B}T_{\rm phys}
 \frac{Z^2\lambda_{\mathrm B}}
 {\sigma(1+\kappa\sigma/2)^2}.
\end{equation}
The paper's reduced temperature is the inverse contact coupling,
\begin{equation}
 T^* \equiv \frac{k_{\mathrm B}T_{\rm phys}}{\epsilon_Y}
 =\frac{\sigma(1+\kappa\sigma/2)^2}{Z^2\lambda_{\mathrm B}}.
 \label{eq:Tstar}
\end{equation}
Thus $\beta\epsilon_Y=1/T^*$.  This is worth remembering: increasing $T^*$
weakens the interaction in thermal units.

With $r^*=r/\sigma$, $r_c^*=r_c/\sigma$, $\kappa^*=\kappa\sigma$, and
$u^*=u/\epsilon_Y$, Eq.~\eqref{eq:paper-potential} becomes
\begin{equation}
u^*_{ij}(r^*)=
\begin{cases}
+\infty,&r^*<1,\\[1mm]
\displaystyle \sgnq\frac{\exp[-\kappa^*(r^*-1)]}{r^*},
&1\le r^*<r_c^*,\\[2mm]
0,&r^*\ge r_c^*.
\end{cases}
\label{eq:reduced-potential}
\end{equation}
In a Metropolis test the exponent is therefore
\begin{equation}
 -\beta\Delta U=-\frac{\Delta U^*}{T^*}.
\end{equation}
The code stores $U^*$ and sets \code{BETA=1/TEMP}; hence the input
\code{TEMP} is $T^*$.

The other reduced variables are
\begin{align}
 \rho^*&=\rho\sigma^3=\frac{N\sigma^3}{V},\\
 p^*&=\frac{p\sigma^3}{\epsilon_Y}
     =\frac{pT^*\sigma^3}{k_{\mathrm B}T_{\rm phys}},\\
 V^*&=\frac{V}{\sigma^3},\qquad L^*=\frac{L}{\sigma}.
\end{align}
Therefore the Fortran input \code{P} is $p^*$ and lengths are measured in
$\sigma$.  For the present runs $\sigma=1$, so the numerical values of $r$
and $r^*$ coincide, but their meanings should not be confused.

\section{Exact map to the Fortran source}
With \code{sig=1}, \code{kappa=6}, and \code{rc=3.5},
Eq.~\eqref{eq:reduced-potential} becomes
\begin{equation}
 u^*_{ij}(r)=\sgnq\frac{e^{-6(r-1)}}{r},
 \qquad 1\le r\le3.5.
 \label{eq:actual-code-potential}
\end{equation}
This is exactly what \code{Source/ener.f} evaluates:
\begin{align}
 \code{pref}&=\exp[-\code{KAPPA}(r-\sigma)],\\
 \code{En}&=\sgnq\,\code{pref}/(r/\sigma).
\end{align}
The branch $r^2<\sigma^2$ returns $10^{20}$, effectively forbidding an
overlap.  The branch $r^2>r_c^2$ returns zero.  The potential is truncated,
not shifted: \code{tailco=.false.} and \code{shift=.false.}.

\subsection{Charge sign, periodicity, and total energy}
\code{Source/eneri.f} chooses $\sgnq=+1$ when \code{PTYPE(i)==PTYPE(j)} and
$-1$ otherwise.  It applies the minimum-image convention independently in
each Cartesian direction.  \code{Source/toterg.f} sums every unordered pair
once, so
\begin{equation}
 U^*=\sum_{i<j}u^*_{ij}(r_{ij}).
\end{equation}
The system contains 500 particles of each type.  For a new lattice the code
assigns alternating IDs to types 1 and 2; restarts store the type explicitly.

\subsection{Pair virial used by the code}
For the finite part of the interaction,
\begin{equation}
 -r\frac{du^*_{ij}}{dr}
 =\sgnq e^{-\kappa(r-\sigma)}\sigma
 \left(\kappa+\frac{1}{r}\right),
\end{equation}
which is the expression in \code{ener.f}.  The reported continuous estimator
is
\begin{equation}
 p^*_{\rm continuous}=\rho^*T^*+\frac{1}{3V^*}
 \sum_{i<j}\left(-r_{ij}\frac{du^*_{ij}}{dr_{ij}}\right).
\end{equation}
It omits the discontinuous hard-core contact contribution and must not be
mistaken for the complete mechanical pressure.  The imposed $p^*$ is still
used correctly in the NpT acceptance rule.

\subsection{Legacy names that do not set the Yukawa strength}
The model line currently reads
\begin{verbatim}
eps sig mass rc kappa
1.0 1.0 1.0 3.5 6.0
\end{verbatim}
The variables \code{eps}, \code{EPS4}, and \code{EPS48} are remnants of the
old Lennard--Jones code.  \code{ENER} does not multiply the Yukawa energy by
\code{eps}.  The Yukawa energy scale is fixed by the reduced-unit definition,
and temperature dependence enters through \code{BETA=1/TEMP}.  In current
runs \code{eps=1} is a harmless placeholder, not an independent coupling.

\section{Packing fraction, density, volume, and box length}
For monodisperse spheres of diameter $\sigma$, the volume of one sphere is
$\pi\sigma^3/6$.  The packing fraction (called $\eta$ here and often $\phi$
elsewhere) is
\begin{equation}
 \boxed{\eta=\frac{N(\pi\sigma^3/6)}{V}
 =\frac{\pi}{6}\rho^*.}
\end{equation}
All required conversions follow:
\begin{align}
 \boxed{\rho^*=\frac{6\eta}{\pi}},\qquad
 \boxed{V^*=\frac{N}{\rho^*}=\frac{N\pi}{6\eta}},\qquad
 \boxed{L^*=\left(\frac{N}{\rho^*}\right)^{1/3}
 =\left(\frac{N\pi}{6\eta}\right)^{1/3}}.
 \label{eq:eta-conversions}
\end{align}
These equations are independent of the Yukawa potential; they are pure
geometry.  Changing $N$ at fixed $\eta$ changes $V$ and $L$, but not $\rho^*$.

\subsection{Worked example: paper state point, $N=1000$ and $\eta=0.526$}
\begin{align}
 \rho^*&=\frac{6(0.526)}{\pi}=1.0045860008,\\
 V^*&=\frac{1000}{1.0045860008}=995.4349346,\\
 L^*&=(995.4349346)^{1/3}=9.98475990.
\end{align}
Thus a newly generated $N=1000$ box at the paper's quoted metastable-liquid
packing fraction should use \code{rho=1.004586}, giving $L\simeq9.98476$.
Using \code{rho=0.526} would be wrong: it would treat packing fraction as
number density and produce $\eta=\pi(0.526)/6\simeq0.2754$.

At the paper's coexistence-liquid value $\eta=0.471$,
\begin{equation}
 \rho^*=0.8995437384,\quad V^*=1111.674683,\quad L^*=10.35919254.
\end{equation}

\subsection{Worked example: current validated $p^*=18$ parent}
The validated restart has $\rho^*=1.03708035$.  Therefore
\begin{equation}
 \eta=\frac{\pi(1.03708035)}{6}=0.543014,qquad
 L^*=\left(\frac{1000}{1.03708035}\right)^{1/3}=9.87936881.
\end{equation}
This state is denser than the paper's $p^*=15$, $\eta\simeq0.526$ liquid.
That is expected because the clean reconstruction currently imposes
$p^*=18$, not $p^*=15$.  Packing fraction is an NpT output at fixed pressure,
not a separately constrained quantity during production.

\section{How initialization actually uses density}
\code{Source/readdat.f} first computes
\begin{equation}
 L^*=\left(\frac{N}{\rho^*_{\rm input}}\right)^{1/3}.
\end{equation}
What happens next depends on \code{ibeg}:
\begin{longtable}{>{\ttfamily}c p{0.73\textwidth}}
\toprule
ibeg & Meaning and effect of the density line\\
\midrule
0 & Generate a new lattice.  The input $\rho^*$ determines the box length.\\
1 & Read a restart and retain its box.  The code replaces the nominal input
density by $N/L_{\rm restart}^3$; the density line does not rescale it.\\
2 & Read a restart and rescale coordinates and box to the requested input
density.  Use deliberately: every distance changes.\\
\bottomrule
\end{longtable}
The clean production workflow uses \code{ibeg=1}, so a line such as
\code{1000 1.0 1.0 18.0 'sc'} does \emph{not} force $\rho^*=1$.  The restart
box is authoritative.  This is why the run output can begin near
$\rho^*=1.037$ even though the nominal line says 1.0.

For NpT production, the target thermodynamic controls are $N$, $T^*$, and
$p^*$.  Density then fluctuates as volume moves are accepted.  An initial
$\eta$ is used only to construct or rescale a starting box; it is not held
fixed by NpT sampling.

\section{NpT volume move in reduced units}
The code proposes a symmetric step in logarithmic volume,
\begin{equation}
 \ln V_n^*=\ln V_o^*+(\xi-1/2)\,\code{Vmax},\qquad 0\le\xi<1,
\end{equation}
then scales all coordinates by $(V_n^*/V_o^*)^{1/3}$.  Its acceptance
exponent is
\begin{equation}
 -\frac{1}{T^*}\left[\Delta U^*+p^*\Delta V^*
ight]
 +(N+1)\ln\left(\frac{V_n^*}{V_o^*}\right).
\end{equation}
The $N+1$ (rather than $N$) is consistent with proposing uniformly in
$\ln V$.  The physical cutoff remains fixed at $3.5\sigma$ during a volume
move; it is not scaled with the box.  A trial is rejected if $L/2\le r_c$.

\section{Current parameter map and audit checklist}
\begin{center}
\begin{tabular}{lll}
\toprule
Physical/reduced quantity & Current value & Code/input location\\
\midrule
$N$ & 1000 & \code{fort.15}, \code{NPART}\\
composition & 500:500 & \code{PTYPE} in restart\\
$\sigma$ & 1 & \code{fort.25}: \code{sig}\\
$T^*$ & 1 & \code{fort.15}: \code{temp}\\
$p^*$ & 18 & \code{fort.15}: \code{p}\\
$\kappa\sigma$ & 6 & \code{fort.25}: \code{kappa}\\
$r_c/\sigma$ & 3.5 & \code{fort.25}: \code{rc}\\
tail correction & off & \code{tailco=.false.}\\
energy shift & off & \code{shift=.false.}\\
\bottomrule
\end{tabular}
\end{center}

Before trusting a run, verify: (i) the restart contains 1000 particles and
500 of each type; (ii) its first line gives the expected $L$; (iii)
$N/L^3$ and $\pi N/(6L^3)$ reproduce the reported $\rho^*$ and $\eta$;
(iv) \code{fort.25} contains $1,1,1,3.5,6$ with shift and tail corrections
off; (v) $L/2>3.5$; and (vi) no pair has $r<1$.

\section{Relation to the LAMMPS Yukawa convention}
LAMMPS writes its plain Yukawa interaction as
\begin{equation}
 U_{ij}^{\rm LAMMPS}(r)=A_{ij}\frac{e^{-\kappa r}}{r}.
\end{equation}
Matching Eq.~\eqref{eq:actual-code-potential} requires
\begin{equation}
 A_{ij}=\sgnq\,\sigma e^{\kappa\sigma}.
\end{equation}
For $\sigma=1$ and $\kappa=6$, $|A|=e^6=403.4287934927351$, positive for
like types and negative for unlike types.  Plain LAMMPS Yukawa has no hard
core, so a cross-check is valid only after confirming that the configuration
has no $r<\sigma$ overlaps.

\section{What was modified relative to the inherited code}
The inherited Lennard--Jones skeleton was adapted so that the working source:
\begin{enumerate}
\item evaluates the signed, contact-referenced Yukawa pair energy in
\code{ener.f} and its analytic pair virial;
\item stores a particle type and assigns equal/opposite signs in
\code{eneri.f};
\item reads $\kappa$ and keeps $r_c$ fixed under volume changes;
\item disables LJ tail corrections and potential shifting for production;
\item preserves types in restart and trajectory output; and
\item records $\eta=\pi\rho^*/6$ alongside density in the newer metrics.
\end{enumerate}
The later clean-replication instrumentation did not rewrite the pair
potential or NpT acceptance physics; it added validation, provenance, and
unambiguous diagnostics.

\section{Reference}
E. Sanz, C. Valeriani, D. Frenkel, and M. Dijkstra, ``Evidence for
Out-of-Equilibrium Crystal Nucleation in Suspensions of Oppositely Charged
Colloids,'' \emph{Physical Review Letters} \textbf{99}, 055501 (2007),
\href{https://doi.org/10.1103/PhysRevLett.99.055501}{doi:10.1103/PhysRevLett.99.055501}.
The local source consulted was \code{Articel-To-Work-With.pdf}, especially
Eq.~(2), Fig.~1, and the surrounding model description.

\end{document}
